% Standalone source-record documentation; no research attempts.
\documentclass[11pt]{article}
\usepackage[margin=25mm]{geometry}
\usepackage{fontspec,unicode-math}
\setmainfont{Latin Modern Roman}
\setmathfont{Latin Modern Math}
\usepackage{amsmath,amssymb,mathtools,xurl,hyperref}
\hypersetup{colorlinks=true,urlcolor=blue,linkcolor=blue}
\setcounter{secnumdepth}{0}
\setlength{\parindent}{0pt}
\setlength{\parskip}{5pt}
\setlength{\emergencystretch}{3em}
\newcommand{\sref}[2]{\hyperlink{source:#1:#2}{[S#2]}}
\newcommand{\cataloguescope}[1]{\paragraph{Source record.}#1}
\newcommand{\floor}[1]{\left\lfloor#1\right\rfloor}
\begin{document}
% UnsolvedMath ID 20001252; source catalogue code AIM-COMPUTATION-0090
\setcounter{section}{20001251}
\section{Degree bounds and the dependence on the modulus}\label{20001252}
{\small\sffamily UnsolvedMath ID 20001252 · Computer Science}
\subsection{Definitions and mathematical statement}

\paragraph{Shared notation.}Standard number systems are $\mathbb N,\mathbb Z,\mathbb Q,\mathbb R,\mathbb C$. Specialized notation and index ranges follow the cited source and the definitions below.


For a prime $p$, $\mathbb F_p[X]$ is the polynomial ring over the field of $p$ elements. In a quotient ring $\mathbb F_p[X]/(h)$, the residue class of $X-s$ is a unit when it is coprime to $h$. These units can generate a finite multiplicative subgroup whose order depends on the modulus and the chosen elements. Group-size lower bounds must state that ambient quotient and the degree and factorization assumptions on $h$.
\paragraph{Record scope and qualifications.}
The catalogue title is retained as metadata. Named certificates, contracts, bounds or frameworks in that title are not asserted as results of the original source extract.
The extraction drops cardinality symbols and does not explicitly give the ambient quotient or necessary modulus hypotheses. The numerical bounds in the historical passage are preserved as source claims rather than asserted for unrestricted $h$. The supplied catalogue marks this record partially solved; its present research status was not independently checked.
The source text below is rendered literally where necessary; damaged formulas and composite extracts are not silently repaired into a new mathematical claim.
\paragraph{Statement recorded in the supplied source.}
Question 2. Consider Fp and d ∈ Z>0, and S ⊂ Fp with S = d. Let h(X) ∈

Fp[X] with h(s) 6 = 0 for all s ∈ S. Let G be the group generated by X − s for s ∈ S. In the original AKS paper, we have G ≥ 2d. There are techniques for getting G ≥ (5.82) d

(using lattice point counting); there seems to be much more room for improvement (perhaps using ABC). (Bernstein, Voloch) Perhaps (log G)/d → ∞

as d → ∞ uniformly over h. \sref{20001252}{2}
\paragraph{Source status.}
The supplied catalogue labels this source record partially solved; that metadata was checked against the supplied snapshot on 2026-10-11. The label is preserved for provenance and is not an independent verification of present open status or mathematical correctness.
\subsection{Short English statement}
The source asks how large a group generated by linear factors can be in an AKS-related setting.
\subsection{Sources}
\begin{itemize}
\item[\hypertarget{source:20001252:1}{S1}] Ulam AI and the UnsolvedMath Contributors, UnsolvedMath: A Curated Collection of Open Mathematics Problems, record 20001252 (AIM-COMPUTATION-0090). Catalogue curation and metadata: CC BY 4.0. The source extract is retained; definitions, English summaries and qualifications are editorial additions.
\url{https://huggingface.co/datasets/ulamai/UnsolvedMath}.
\item[\hypertarget{source:20001252:2}{S2}] AIM Workshop Problem Lists, source linked by catalogue record AIM-COMPUTATION-0090.
\url{https://aimath.org/WWN/primesinp/primesinp.pdf}.
\end{itemize}
\paragraph{Preparation and review.}
The original source and UnsolvedMath contributors retain attribution for the source extract. Codex (GPT-6; exact serving revision not exposed) assisted with definitions, English summaries, scope notes and formatting. No new human mathematical authorship or independently verified solution is asserted. Permanent source ID 20001252 and code AIM-COMPUTATION-0090 checked against the complete supplied registry and current repository. Distinct upstream records, including repeated questions and broad or damaged extracts, are retained for complete coverage of the source CS category. A repeated record is not a new independent conjecture. Complete mathematical-equivalence and current-literature review has not been performed. Preparation and final review follow contribution prompts C and D, with source gaps disclosed. Structural and rendering checks do not establish a new theorem or open status.
\label{end:20001252}
\end{document}
